\documentclass[UTF8,12pt]{ctexart}
\usepackage[paperwidth=240mm,paperheight=200mm,margin=14mm,top=8mm]{geometry}
\usepackage{amsmath,amssymb,mathtools}
\usepackage{xcolor}
\pagestyle{empty}
\setlength{\parindent}{0pt}
\setlength{\parskip}{0pt}
\newcommand{\qn}[1]{\textbf{\large #1}\ }
\xeCJKDeclareCharClass{CJK}{"2460 -> "24FF}
\begin{document}
\qn{1990.1}
曲线 $\begin{cases} x=\cos^3 t, \\ y=\sin^3 t \end{cases}$ 上对应于 $t=\dfrac{\pi}{6}$ 处的法线方程是______.

\vfill
\newpage
\qn{1990.2}
设 $y=\mathrm{e}^{\tan\frac{1}{x}}\sin\dfrac{1}{x}$，则 $y'=$______.

\vfill
\newpage
\qn{1990.3}
$\int_0^1 x\sqrt{1-x}\,\mathrm{d}x=$______.

\vfill
\newpage
\qn{1990.4}
下列两个积分的大小关系是：$\int_{-2}^{-1}\mathrm{e}^{-x^3}\,\mathrm{d}x$______$\int_{-2}^{-1}\mathrm{e}^{x^3}\,\mathrm{d}x$.

\vfill
\newpage
\qn{1990.5}
设函数 $f(x)=\begin{cases} 1, & |x|\leqslant 1, \\ 0, & |x|>1, \end{cases}$ 则函数 $f[f(x)]=$______.

\vfill
\newpage
\qn{1990.6}
已知 $\lim\limits_{x\to\infty}\left(\dfrac{x^2}{x+1}-ax-b\right)=0$，其中 $a,b$ 是常数，则（　　）
\\
\noindent (A) $a=1,\,b=1$.　　(B) $a=-1,\,b=1$.
\\
\noindent (C) $a=1,\,b=-1$.　　(D) $a=-1,\,b=-1$.

\vfill
\newpage
\qn{1990.7}
设函数 $f(x)$ 在 $(-\infty,+\infty)$ 上连续，则 $\mathrm{d}\left[\int f(x)\,\mathrm{d}x\right]$ 等于（　　）
\\
\noindent (A) $f(x)$.　　(B) $f(x)\,\mathrm{d}x$.
\\
\noindent (C) $f(x)+C$.　　(D) $f'(x)\,\mathrm{d}x$.

\vfill
\newpage
\qn{1990.8}
已知函数 $f(x)$ 具有任意阶导数，且 $f'(x)=[f(x)]^2$，则当 $n$ 为大于 2 的正整数时，$f(x)$ 的 $n$ 阶导数 $f^{(n)}(x)$ 是（　　）
\\
\noindent (A) $n![f(x)]^{n+1}$.　　(B) $n[f(x)]^{n+1}$.
\\
\noindent (C) $[f(x)]^{2n}$.　　(D) $n![f(x)]^{2n}$.

\vfill
\newpage
\qn{1990.9}
设 $f(x)$ 是连续函数，且 $F(x)=\int_x^{\mathrm{e}^{-x}}f(t)\,\mathrm{d}t$，则 $F'(x)$ 等于（　　）
\\
\noindent (A) $-\mathrm{e}^{-x}f(\mathrm{e}^{-x})-f(x)$.　　(B) $-\mathrm{e}^{-x}f(\mathrm{e}^{-x})+f(x)$.
\\
\noindent (C) $\mathrm{e}^{-x}f(\mathrm{e}^{-x})-f(x)$.　　(D) $\mathrm{e}^{-x}f(\mathrm{e}^{-x})+f(x)$.

\vfill
\newpage
\qn{1990.10}
设 $F(x)=\begin{cases} \dfrac{f(x)}{x}, & x\neq 0, \\ f(0), & x=0, \end{cases}$ 其中 $f(x)$ 在 $x=0$ 处可导，$f'(0)\neq 0$，$f(0)=0$，则 $x=0$ 是 $F(x)$ 的（　　）
\\
\noindent (A) 连续点.　　(B) 第一类间断点.
\\
\noindent (C) 第二类间断点.　　(D) 连续点或间断点不能由此确定.

\vfill
\newpage
\qn{1990.11}
已知 $\lim\limits_{x\to\infty}\left(\dfrac{x+a}{x-a}\right)^x=9$，求常数 $a$.

\vfill
\newpage
\qn{1990.12}
求由方程 $2y-x=(x-y)\ln(x-y)$ 所确定的函数 $y=y(x)$ 的微分 $\mathrm{d}y$.

\vfill
\newpage
\qn{1990.13}
求曲线 $y=\dfrac{1}{1+x^2}$（$x>0$）的拐点.

\vfill
\newpage
\qn{1990.14}
计算 $\int\dfrac{\ln x}{(1-x)^2}\,\mathrm{d}x$.

\vfill
\newpage
\qn{1990.15}
求微分方程 $x\ln x\,\mathrm{d}y+(y-\ln x)\,\mathrm{d}x=0$ 满足条件 $\left.y\right|_{x=\mathrm{e}}=1$ 的特解.
在椭圆 $\dfrac{x^2}{a^2}+\dfrac{y^2}{b^2}=1$ 的第一象限部分上求一点 $P$，使该点处的切线，椭圆及两坐标轴所围图形面积为最小（其中 $a>0,\,b>0$）.
证明：当 $x>0$ 时，有不等式 $\arctan x+\dfrac{1}{x}>\dfrac{\pi}{2}$.
设 $f(x)=\int_1^x\dfrac{\ln t}{1+t}\,\mathrm{d}t$，其中 $x>0$，求 $f(x)+f\left(\dfrac{1}{x}\right)$.
过点 $P(1,0)$ 作抛物线 $y=\sqrt{x-2}$ 的切线，该切线与上述抛物线及 $x$ 轴围成一平面图形，求此平面图形绕 $x$ 轴旋转一周所成旋转体的体积.
求微分方程 $y''+4y'+4y=\mathrm{e}^{ax}$ 的通解，其中 $a$ 为实数.

\vfill
\newpage
\end{document}